
\Slide{1}{本讲的任务}
不动点把“选择相互协调”写成一个数学条件。先区分函数的不动点等式和对应的不动点包含，再逐项核对非空、紧、凸、连续性。

\Slide{2}{依赖的知识}
集合、连续性、紧性与凸性缺一不可。高维 Brouwer 是比介值定理更深的定理，课件只陈述，本讲补充证明路线及 Kakutani 的推导。

\Slide{3}{函数不动点}
$x=f(x)$ 意味输入与输出同一个对象。例 $f(x)=1-x$ 在 $[0,1]$ 的不动点是 $1/2$；而 $f(x)=x+1$ 在 $\R$ 无不动点。不是每个映射都有不动点。

\Slide{4}{Brouwer 定理的条件}
非空紧凸 $S\subset\R^N$ 上的连续自映射必有不动点。各条件要检查：$f(S)\subseteq S$ 是自映射条件；在球面旋转可无不动点，说明凸性不能仅用闭有界代替。附录给单纯形上的组合证明。

\Slide{5}{球与一般紧凸集的归约纠错}
一般 $S$ 可能是线段或单点，不能无条件同胚于 $N$ 维球。正确归约用投影：设 $D$ 为包含 $S$ 的大闭球，$P:D\to S$ 为最近点投影。最近点存在由紧性；唯一性由平方距离的中点恒等式，否则两个不同最优点中点更近。最优性给 $(x-Px)\cdot(z-Px)\le0$。对 $x,y$ 分别代入 $Py,Px$ 并相加得 $\norm{Px-Py}^2\le(x-y)\cdot(Px-Py)$，Cauchy--Schwarz 给 $\norm{Px-Py}\le\norm{x-y}$，故投影连续。对 $f\circ P:D\to S\subset D$ 用球版定理，若 $x=f(Px)$ 则 $x\in S,Px=x$，即 $x=f(x)$。

\Slide{6}{一维证明与值域纠错}
设 $h:[-1,1]\to[-1,1]$ 连续，$g(x)=h(x)-x$ 为实值连续函数，值域不一定在 $[-1,1]$，原页此处应改为 $\R$。端点有 $g(-1)\ge0,g(1)\le0$，若等于0已得到不动点，否则介值定理给内点零点。$h'$ 在课件只是新函数符号，不能误读为导数。

\Slide{7}{对应的不动点}
对应 $\Phi:S\rightrightarrows S$ 的不动点是 $x\in\Phi(x)$，不要求 $\Phi(x)=\{x\}$。若值集是最优响应，含义是当前选择本身就是允许的最优响应之一。

\Slide{8}{Kakutani 与完整推导}
采用标准安全版本：$S$ 非空紧凸，$\Phi$ 非空紧凸值且上半连续，则有不动点。闭值不能遗漏；在紧陪域中闭图与闭值上半连续相互联系。证明：对每个 $n$ 取有限 $1/n$ 网 $a_i\in S$，每个选 $b_i\in\Phi(a_i)$。设 $w_i(x)=\max(0,2/n-\norm{x-a_i})$，分母 $\sum_iw_i(x)>0$，定义连续 $f_n(x)=\sum_iw_i(x)b_i/\sum_iw_i(x)\in S$。Brouwer 给 $x_n=f_n(x_n)$。抽子列 $x_n\to x$。对任意 $\eps>0$，上半连续使所有充分接近 $x$ 的输入 $a$ 满足 $\Phi(a)\subset\Phi(x)+B(0,\eps)$。参与正权重的中心满足 $\norm{a_i-x_n}<2/n$，所以最终均足够接近 $x$。右侧为凸集（值凸与球凸），故其凸组合 $x_n$ 也在其中；于是 $\dist(x_n,\Phi(x))<\eps$。取极限及任意 $\eps$，闭值性给 $x\in\Phi(x)$。注意这里不能从“闭图”直接推出任意移动凸组合仍归于原值，真正使用的是上半连续、凸性与闭性。

\clearpage\section*{进阶附录：Sperner 到 Brouwer 的证明}
令标准单纯形 $\Delta^n=\{x_i\ge0,\sum_{i=0}^n x_i=1\}$。把它作有限相容三角剖分，每顶点标号 $0,\ldots,n$，规定若顶点的第 $i$ 坐标为0便不能标 $i$。
\paragraph{Sperner 引理及计数证明} 完全标号（包含全部 $n+1$ 种标号）的小 $n$ 单纯形个数为奇数。对维数归纳，维数0显然。统计每个小 $n$ 单纯形具有标号恰为 $0,\ldots,n-1$ 的面数。完全标号者有1个；缺少标号 $n$ 而具有其他全部标号者恰有2个，因为其中某一标号重复；其他有0个。故总数模2等于完全标号者个数模2。内面由相邻两个单纯形各计一次，模2消去。边界合规面只能在大单纯形的 $x_n=0$ 面上，因为其他面缺少某个必需标号；归纳假设给其数量为奇数，得证。
\paragraph{从引理构造不动点} 对连续 $f:\Delta^n\to\Delta^n$，若剖分某顶点已满足 $f(v)=v$ 就完成。否则某坐标 $i$ 满足 $v_i>f_i(v)$（两者坐标和相同），给该点选一个这样的标号。因 $v_i=0$ 不可能大于非负 $f_i(v)$，边界规则成立。取网格直径趋于0的剖分（反复重心剖分使最大直径至多逐次乘 $n/(n+1)$），由引理每层有完全标号小单纯形。紧性抽子列使其0号顶点趋于 $x$；直径趋0使其余各号顶点也趋于同一个 $x$。对标号 $i$ 的顶点有 $v_i>f_i(v)$，极限给 $x_i\ge f_i(x)$。全部坐标和相同，故每项只能相等，$x=f(x)$。
\paragraph{从单纯形到一般集合} 给紧凸 $S\subset\R^N$ 取包含它的大 $N$ 单纯形 $T$，与标准单纯形有仿射同构。第5页的最近点投影 $P:T\to S$ 连续，对 $f\circ P:T\to T$ 用已证单纯形版，得到的固定点自动在 $S$，于是为 $f$ 的固定点。这样不需假设低维 $S$ 与全维球同胚。
\paragraph{考试自测} (1) 若 $S=(0,1)$，$f(x)=(x+1)/2$，为何失败？解：不紧，不动点方程只解得不在域中的1。(2) 唯一性是否成立？解：恒等函数的每一点都固定。(3) Kakutani 证明哪步用凸值？解：把附近选值的凸组合仍放入 $\Phi(x)+B(0,\eps)$；哪步用闭值？最后由距离为0得到属于值集。
